% GENERATED FILE. Edit canonical lessons, then run npm run generate:books.
% canonical-source-hash: fnv1a64:c5e26abc5e2ed2ac
% canonical-lessons: JA-C276-ouame, JA-C276-oumizu, JA-C276-afureru, JA-C276-gake, JA-C276-taihen

\chapter{The River Rose}
\label{ch:ja-the-river-rose}
\hlchaptermodality{276}

\begin{chapteropening}
\textbf{By the end of this chapter:} \emph{I can say what heavy rain did to the river and the hills.}

Complete the last lesson of chapter 276: I can say what heavy rain did to the river and the hills.
\end{chapteropening}

\section[\texorpdfstring{\ja{おおあめ} (ōame)}{おおあめ (ōame)}]{\ja{おおあめ} (ōame) — heavy rain; a downpour}

\label{lesson:JA-C276-ouame}



\begin{quote}
Before the new one: say the Japanese for to float, then the Japanese for once more.
\end{quote}

\subsection*{You'll want to know: \ja{おおあめ}}
\textbf{\ja{おおあめ}} — \emph{ōame} — \textquotedblleft{}heavy rain; a downpour\textquotedblright{}.

\begin{quote}
\textbf{\ja{きのうは} \ja{おおあめでした。}} — \emph{kinō wa ōame deshita} — \textquotedblleft{}It poured yesterday.\textquotedblright{}
\end{quote}

\begin{cousinweb}
\textbf{\ja{おおあめ}} (\emph{ōame}) is a native Japanese word (wago), inherited rather than borrowed from Chinese, built from \emph{ō}, \textquotedblleft{}big\textquotedblright{}, and \emph{ame}, \textquotedblleft{}rain\textquotedblright{} (JMdict, KANJIDIC2).
\end{cousinweb}

\begin{grammarlens}[title={What you've built}]
One more word for telling what happened.
\end{grammarlens}

\subsection*{Your turn}
\begin{itemize}
\raggedright
  \item \emph{Say it:} \emph{ōame}
  \item \emph{Say it:} \emph{ōame}, once more
  \item \emph{Say it:} \emph{kinō wa ōame deshita} — It poured yesterday.
  \item \emph{Recall:} say the Japanese for to float, then the Japanese for once more, then say \emph{ōame} again
\end{itemize}

\begin{tcolorbox}[breakable,colback=teal!4,colframe=teal!35!black,title={Before you move on}]
What does \ja{おおあめ} mean? (\textquotedblleft{}Heavy rain; a downpour\textquotedblright{}.) Say it once more.
\end{tcolorbox}
\section[\texorpdfstring{\ja{おおみず} (ōmizu)}{おおみず (ōmizu)}]{\ja{おおみず} (ōmizu) — a flood}

\label{lesson:JA-C276-oumizu}



\begin{quote}
Before the new one: say the Japanese for once more, then the Japanese for heavy rain.
\end{quote}

\subsection*{You'll want to know: \ja{おおみず}}
\textbf{\ja{おおみず}} — \emph{ōmizu} — \textquotedblleft{}a flood\textquotedblright{}.

\begin{quote}
\textbf{\ja{おおみずの} \ja{とき、}\ja{日本人}\ja{の} \ja{ともだちが} \ja{たすけに} \ja{きました。}} — \emph{ōmizu no toki, nihonjin no tomodachi ga tasuke ni kimashita} — \textquotedblleft{}In the flood, a Japanese friend came to help.\textquotedblright{}
\end{quote}

\begin{cousinweb}
\textbf{\ja{おおみず}} (\emph{ōmizu}) is a native Japanese word (wago), inherited rather than borrowed from Chinese, built from \emph{ō}, \textquotedblleft{}big\textquotedblright{}, and \emph{mizu}, \textquotedblleft{}water\textquotedblright{} (JMdict, KANJIDIC2).
\end{cousinweb}

\begin{grammarlens}[title={What you've built}]
One more word for telling what happened.
\end{grammarlens}

\subsection*{Your turn}
\begin{itemize}
\raggedright
  \item \emph{Say it:} \emph{ōmizu}
  \item \emph{Say it:} \emph{ōmizu}, once more
  \item \emph{Say it:} \emph{ōmizu no toki, nihonjin no tomodachi ga tasuke ni kimashita} — In the flood, a Japanese friend came to help.
  \item \emph{Recall:} say the Japanese for once more, then the Japanese for heavy rain, then say \emph{ōmizu} again
\end{itemize}

\begin{tcolorbox}[breakable,colback=teal!4,colframe=teal!35!black,title={Before you move on}]
What does \ja{おおみず} mean? (\textquotedblleft{}A flood\textquotedblright{}.) Say it once more.
\end{tcolorbox}
\section[\texorpdfstring{\ja{あふれる} (afureru)}{あふれる (afureru)}]{\ja{あふれる} (afureru) — to overflow}

\label{lesson:JA-C276-afureru}



\begin{quote}
Before the new one: say the Japanese for heavy rain, then the Japanese for a flood.
\end{quote}

\subsection*{You'll want to know: \ja{あふれる}}
\textbf{\ja{あふれる}} — \emph{afureru} — \textquotedblleft{}to overflow\textquotedblright{}.

\begin{quote}
\textbf{\ja{いけの} \ja{みずが} \ja{あふれています。}} — \emph{ike no mizu ga afurete imasu} — \textquotedblleft{}The pond is overflowing.\textquotedblright{}
\end{quote}

Polite \textbf{\ja{あふれます}}, \emph{afuremasu}; an ichidan verb. Past (polite): \textbf{\ja{かわが} \ja{あふれました。}}, \emph{kawa ga afuremashita}.

\begin{cousinweb}
\textbf{\ja{あふれる}} (\emph{afureru}) is a native Japanese word (wago), inherited rather than borrowed from Chinese; its kanji was attached later as a kun'yomi (JMdict).
\end{cousinweb}

\begin{grammarlens}[title={What you've built}]
One more word for telling what happened.
\end{grammarlens}

\subsection*{Your turn}
\begin{itemize}
\raggedright
  \item \emph{Say it:} \emph{afureru}
  \item \emph{Say it:} \emph{afureru}, once more
  \item \emph{Say it:} \emph{ike no mizu ga afurete imasu} — The pond is overflowing.
  \item \emph{Recall:} say the Japanese for heavy rain, then the Japanese for a flood, then say \emph{afureru} again
\end{itemize}

\begin{tcolorbox}[breakable,colback=teal!4,colframe=teal!35!black,title={Before you move on}]
What does \ja{あふれる} mean? (\textquotedblleft{}To overflow\textquotedblright{}.) Say it once more.
\end{tcolorbox}
\section[\texorpdfstring{\ja{がけ} (gake)}{がけ (gake)}]{\ja{がけ} (gake) — a cliff; a steep slope}

\label{lesson:JA-C276-gake}



\begin{quote}
Before the new one: say the Japanese for a flood, then the Japanese for to overflow.
\end{quote}

\subsection*{You'll want to know: \ja{がけ}}
\textbf{\ja{がけ}} — \emph{gake} — \textquotedblleft{}a cliff; a steep slope\textquotedblright{}.

\begin{quote}
\textbf{\ja{おおあめの} \ja{あと、がけが} \ja{くずれました。}} — \emph{ōame no ato, gake ga kuzuremashita} — \textquotedblleft{}After the downpour, the slope gave way.\textquotedblright{}
\end{quote}

\begin{cousinweb}
\textbf{\ja{がけ}} (\emph{gake}) is a native Japanese word (wago), inherited rather than borrowed from Chinese; its kanji was attached later as a kun'yomi (JMdict).
\end{cousinweb}

\begin{grammarlens}[title={What you've built}]
One more word for telling what happened.
\end{grammarlens}

\subsection*{Your turn}
\begin{itemize}
\raggedright
  \item \emph{Say it:} \emph{gake}
  \item \emph{Say it:} \emph{gake}, once more
  \item \emph{Say it:} \emph{ōame no ato, gake ga kuzuremashita} — After the downpour, the slope gave way.
  \item \emph{Recall:} say the Japanese for a flood, then the Japanese for to overflow, then say \emph{gake} again
\end{itemize}

\begin{tcolorbox}[breakable,colback=teal!4,colframe=teal!35!black,title={Before you move on}]
What does \ja{がけ} mean? (\textquotedblleft{}A cliff; a steep slope\textquotedblright{}.) Say it once more.
\end{tcolorbox}
\section[\texorpdfstring{\ja{たいへん} (taihen)}{たいへん (taihen)}]{\ja{たいへん} (taihen) — awful; serious}

\label{lesson:JA-C276-taihen}



\begin{quote}
Before the new one: say the Japanese for to overflow, then the Japanese for a cliff.
\end{quote}

\subsection*{You'll want to know: \ja{たいへん}}
\textbf{\ja{たいへん}} — \emph{taihen} — \textquotedblleft{}awful; serious\textquotedblright{}.

\begin{quote}
\textbf{\ja{おおあめで、きのうは} \ja{たいへんでした。}} — \emph{ōame de, kinō wa taihen deshita} — \textquotedblleft{}With the downpour, yesterday was awful.\textquotedblright{}
\end{quote}

A na-adjective: \ja{な} before a noun.

\begin{cousinweb}
\textbf{\ja{たいへん}} (\emph{taihen}) is Sino-Japanese (kango): two Chinese characters read with their on'yomi, readings borrowed from Chinese. They are \emph{tai}, \textquotedblleft{}big\textquotedblright{} (Middle Chinese \emph{dajH}) and \emph{hen}, \textquotedblleft{}unusual\textquotedblright{} (\emph{pjenH}) (KANJIDIC2; Middle Chinese in Baxter's transcription).
\end{cousinweb}

\begin{grammarlens}[title={What you've built}]
One more word for telling what happened.
\end{grammarlens}

\subsection*{Your turn}
\begin{itemize}
\raggedright
  \item \emph{Say it:} \emph{taihen}
  \item \emph{Say it:} \emph{taihen}, once more
  \item \emph{Say it:} \emph{ōame de, kinō wa taihen deshita} — With the downpour, yesterday was awful.
  \item \emph{Recall:} say the Japanese for to overflow, then the Japanese for a cliff, then say \emph{taihen} again
\end{itemize}

\begin{tcolorbox}[breakable,colback=teal!4,colframe=teal!35!black,title={Before you move on}]
What does \ja{たいへん} mean? (\textquotedblleft{}Awful; serious\textquotedblright{}.) Say it once more.
\end{tcolorbox}
